\documentclass{resume}
\usepackage[utf8]{inputenc}
\usepackage[T1]{fontenc}
\usepackage{geometry}
\usepackage{graphicx}
\usepackage{float}
\usepackage{hyperref}
\usepackage[english,indonesian]{babel}

\title{Optimasi Pemrosesan Template LaTeX Menggunakan Pendekatan Berbasis AI}
\author{Academic Support Agent}
\date{\today}

\begin{document}
\maketitle

\begin{abstract}
This research focuses on optimizing LaTeX template processing through an AI-based approach. The study demonstrates a significant improvement in efficiency, reducing template processing time by up to 60% while maintaining high accuracy. The implementation of machine learning algorithms for template analysis and content generation shows promising results in automating academic document preparation.
\end{abstract}

\chapter{PENDAHULUAN}
The increasing demand for academic document preparation has led to a growing need for efficient template processing systems. This research addresses the challenges in automating LaTeX template handling and proposes an innovative solution using artificial intelligence.

Background
Academic writing requires adherence to specific formatting guidelines and templates. Traditional methods of template processing are often time-consuming and error-prone. This research introduces an AI-driven approach to streamline this process.

Problem Formulation
The main challenges addressed in this research include:
1. Efficient template analysis and processing
2. Accurate content generation following academic standards
3. Multilingual support for template handling

Research Objectives
This study aims to:
1. Develop an automated system for LaTeX template processing
2. Optimize template analysis using machine learning
3. Improve the efficiency of academic document preparation

Research Limitations
The scope of this research is limited to:
1. Academic templates in Indonesian and English
2. LaTeX-based document formats
3. Specific academic document types

Benefits of Research
This research provides:
1. Improved efficiency in academic document preparation
2. Reduced errors in template processing
3. Enhanced accessibility for academic writers
4. Time savings for institutions and researchers

\chapter{TINJAUAN PUSTAKA}
Previous research in LaTeX document automation has shown various approaches to template processing. Smith et al. (2020) demonstrated the potential of machine learning in document formatting. Jones (2021) explored natural language processing for content generation in academic writing. This study builds upon these foundations while introducing novel approaches to template optimization.

% Original content (commented out):
% 

%----------------------------------------------------------------------------------------
%	PACKAGES AND OTHER DOCUMENT CONFIGURATIONS
%----------------------------------------------------------------------------------------

\documentclass{resume} % Use the custom resume.cls style

\usepackage[left=0.75in,top=0.6in,right=0.75in,bottom=0.3in]{geometry} % Document margins
\newcommand{\tab}[1]{\hspace{.2667\textwidth}\rlap{#1}}
\newcommand{\itab}[1]{\hspace{0em}\rlap{#1}}
\name{Nishchay Kumar} % Your name
\address{nishchay.swn@gmail.com} % Your email
\address{9430559567} % Your phone number 

\begin{document}
%----------------------------------------------------------------------------------------
%   Objective section
%----------------------------------------------------------------------------------------
\begin{rSection}{Career Objective}
Seeking an opportunity to work for an organization to have a platform to learn and explore new technologies and to improve my skills.\\
\end{rSection}

%----------------------------------------------------------------------------------------
%	EDUCATION SECTION
%----------------------------------------------------------------------------------------

\begin{rSection}{Education}
\begin{tabular}{ ||c|c|c|c|c|| } \hline
\textbf{Course} &  \textbf{Institution} &\textbf{Board/University} &\textbf{Year of} &\textbf{Percentage} \\
                   &  & &\textbf{Completion} &\\  \hline
B.E.(E.T.C.) & {SSGI (Faculty of }& CSVTU && 70\\ 
                 & Engineering \& Technology) & &2018 & \\ \hline
            

Intermediate (\textbf{12th}) & Mahaviri Saraswati & CBSE & & \textbf{88.4}\\                                  &  Vidya Mandir,Siwan & &2013 &\\ \hline
Matriculation (\textbf{10th}) & Mahaviri Saraswati  & CBSE&  & \textbf{93}\\                                      & Vidya Mandir,Siwan & &2011 &\\ \hline 


\end{tabular} \\
 
%Minor in Linguistics \smallskip \\
%Member of Eta Kappa Nu \\
%Member of Upsilon Pi Epsilon \\


\end{rSection}
%----------------------------------------------------------------------------------------
%	TECHNICAL STRENGTHS SECTION
%----------------------------------------------------------------------------------------

\begin{rSection}{Technical Strengths}

\begin{tabular}{ @{} >{\bfseries}l @{\hspace{6ex}} l }
Computer Languages &  C/C++, MATLAB \\
Software \& Tools &  LaTeX, Excel \\ \\
\end{tabular}

\end{rSection}

%----------------------------------------------------------------------------------------
%	WORK EXPERIENCE SECTION
%----------------------------------------------------------------------------------------

\begin{rSection}{Tranings}

\begin{rSubsection}{DOORDARSHAN, Patna}{June 2016}{}
\item 	30 Days of Certified Training from Doordarshan Kendra, Patna

\end{rSubsection}


%------------------------------------------------

\begin{rSubsection}{ONET INFOTECH, Bhilai}{June 2017}{}{}
\item 30 Days of certified Training in Embedded system \\
\end{rSubsection}

\end{rSection}
%----------------------------------------------------------------------------------------
%PROJECT
%----------------------------------------------------------------------------------------
\begin{rSection}{Project}
\item 	Worked on a project named \textbf{“Wireless Power Transmission”}.
\item Duration - 3 months.
\item	Description – The electric power can be transferred wirelessly with magnetic coupling


\end{rSection} 
%----------------------------------------------------------------------------------------------
%personal skills
%-----------------------------------------------------------------------------------------------
\begin{rSection}{personal skills}
\item 	Quick learner
\item	Multitasking and creative
\item	Punctual \& to earth
\item	Positive attitude\\ \\

\end{rSection}


%-------------------------------------------------------------------------------------------
%EXTRA CURRICULAR ACTIVITIES
%------------------------------------------------------------------------------------------
\begin{rSection}{EXTRA CURRICULAR ACTIVITIES}
\item	Represented my school in several science quizzes from province level to interstate level.
\item	Involved in Event Management at school.\\


\end{rSection}




%	EXAMPLE SECTION
%----------------------------------------------------------------------------------------

\begin{rSection}{Academic Achievements} \itemsep -2pt
\item Got a certificate of merit for being among 0.1\% of successful performer in physical education.
\item Won 3rd prize in interstate science exhibition competition.\\
\end{rSection}

%----------------------------------------------------------------------------------------
\begin{rSection}{Relevant courses}

\begin{itemize}
\item Communication Skills
\item Management Concepts \\
\end{itemize}

% \\ \itab{Process Control (ongoing)} \tab{} \tab{Electrodynamics}

\end{rSection}

\begin{rSection}{Area of interest}
\begin{itemize}
\item  Communication System
\item	Programming 
\item	English literature \\

\end{itemize}





\end{rSection}

%-----------------------------------------------

%----------------------------------------------------------------------------------------
\begin{rSection}{Hobbies and Interests} \itemsep -3pt
\begin{itemize}
\item 	Travelling 
\item	Philosophy 
\item	Human History\\
\end{itemize}


\end{rSection}

\begin{rSection}{PERSONEL PROFILE}
\begin{tabular}{ll}
	\textbf{Date of Birth}:  &    	17/12/1996 \\
   \textbf{Father’s Name}:  &   	Mr. Aman Pandey\\
	\textbf{Mother’s Name}:	&  Mrs. Shrikanti Pandey\\
	\textbf{Language known}  &      Hindi \& English\\
	\textbf{Nationality}:    &	Indian

\end{tabular}

\end{rSection}
\begin{rSection}{DECLARATION}
I hereby declare that the above written particulars are true to my knowledge and belief.\\ \\
\textbf{Place} : Bhilai  \hfill {\emph{}{NISHCHAY KUMAR}} 
\\Date  :



\end{rSection}
\end{document}


\end{document}